\section{Discussion}
\label{sec:discussion}

\subsection{Why simple adaptation wins in this setting}

The evidence favours task adaptation, not a general rejection of
Transformers. The controlled distribution is narrow, the training split
contains 160 rollouts and 4,036 usable windows, the pretrained MultiPath
representation is strong, and the explicit history covers only one second. A
retrained output head can therefore correct substantial domain mismatch
without learning a new interaction representation. The sequence ablations
show that \Tone{} and \Ttwo{} do use their added input; their lower rank is not
explained by complete input neglect. Conversely, the B1 shuffles indicate that
its aggregate gain is raster-dominant rather than evidence of effective target
history modelling.

This interpretation is compatible with two strands of prior evidence. Large
interaction Transformers obtain their gains with richer multi-agent/map inputs
and much larger benchmark distributions
\citep{yuan2021agentformer,ngiam2022scene,nayakanti2023wayformer,shi2022mtr},
whereas cross-dataset analysis shows that domain transfer itself can cause a
large accuracy loss \citep{feng2024unitraj}. The present result adds the
boundary that, for a narrow in-distribution adaptation with a fixed MultiPath
interface, correcting the final head can be more valuable than adding a small
sequence adapter.

The comparison also has an important fairness boundary. The MLP/Transformer
pairs approximately match adapter capacity within two output scopes, but B1
has substantially more trainable parameters and 10/15 runs select a checkpoint
at the epoch ceiling. The result therefore supports simple adaptation as the
best tested configuration under the frozen budget, not an architecture-level
causal claim.

\subsection{Why offline gain does not imply uniform closed-loop gain}

R3 first establishes that the deployed manipulation is real: B1 improves all
40/40 in-loop prediction checks. The failure of uniform H3 translation cannot
therefore be dismissed as a missing predictor difference. The planner consumes
a calibrated distribution rather than ADE alone; risk allocation transforms
that distribution into constraints, controller acceptance/fallback and active
constraints limit candidate motion, and the supervisor can change execution. Because B1 and B0 also use
different frozen calibration, the valid estimand is the deployed predictor
stack rather than neural weights alone.

The result agrees with the task-relevance argument that only some forecast
errors alter downstream cost \citep{farid2023task,nakamura2025regret} and with
closed-loop benchmarks that reject open-loop distance as a sufficient planning
metric \citep{caesar2021nuplan}. It also independently reproduces the direction
of the recent, broader open-loop/closed-loop mismatch reported by
\citet{bouzidi2025closing}. Unlike integrated end-to-end training such as DIPP
\citep{huang2023dipp}, however, this experiment preserves the modular predictor
interface and attributes the observed effect only to the frozen stack-level
treatment.

The corrected result is deliberately modest: two of eight cells move both
completion and separation in the preferred direction, and the pattern depends
on policy and target style. Legacy response-tail and timing analyses supply
plausible mechanisms---calibration failure and simultaneous changes in solver,
supervisor and physical margin---but are not pooled with R3 and do not identify
a single causal component.

\subsection{Adaptive risk remains useful without dominance}

H4's negative universal result does not make adaptive risk valueless. It
dominates a fixed comparator in three prespecified contexts and fails to do so
in nine. Its defensible interpretation is therefore a context-dependent
operating point. Comparing it with aggressive, medium and conservative fixed
settings prevents a favourable single comparator from being presented as
general superiority. The result directly refines the project's original
planning hypothesis while retaining a substantive role for adaptive policy.

This is consistent with risk allocation as a mechanism for reducing unnecessary
conservatism rather than a theorem of universal empirical dominance
\citep{ono2008risk,nair2025predictive}. Formal SMPC safety and feasibility
results depend on their modelling assumptions \citep{ren2025recursively}; the
finite CARLA comparison instead asks which executed operating point is preferred
under the shared corrected controller and supervisor.

\subsection{What corrected-supervisor application authority explains}

% Provisional tracked insertion. The final-only evidence builder atomically
% replaces this file after every scientific and provenance gate passes.
\subsection{Supervisor-authority interpretation pending scientific closure}

The bounded interpretation of supervisor authority is intentionally unavailable
in this non-submission draft. This placeholder contains no outcome estimate,
count, percentage, probability value, effect direction, ranking or superiority
claim. The release audit remains fail-closed until the preregistered authority
matrix, source hashes and paper-integration build have all passed.


\subsection{Threats to validity}

\begin{table}[t]
  \caption{Threats, mitigations and remaining boundaries}
  \label{tab:threats}
  \centering
  \small
  \begin{tabular}{@{}L{0.25\linewidth}L{0.31\linewidth}L{0.34\linewidth}@{}}
    \toprule
    Threat & Mitigation & Remaining boundary \\
    \midrule
    315 overlapping offline windows & Rollout-macro points and five init-level paired effects & Effective inferential unit is five test init groups; minimum exact $p=0.0625$ \\
    Five corrected-R3 init groups & Paired effects, cluster intervals and exact tests & Minimum two-sided exact $p=0.0625$ \\
    Single map and junction & Predictor, policy and target-style factorial & No cross-map or real-road claim \\
    Configuration fairness & Two matched adapter scopes and full parameter/epoch audit & No pure architecture causality \\
    Predictor calibration differs & Stack-level estimand declared before R3 & No weight-only B1 effect \\
    Coupled solver/supervisor & Preregistered corrected-supervisor application-authority on/off SF4 intervention & The DID covers adaptive versus the criticised fixed-medium pair only; estimator-to-risk path and SMPC constraints remain in both arms \\
    Solver bypass and acceptance & Raw classification; factual attempts alone define finite CasADi timing and controller acceptance; bypass and raw status are separate & Acceptance is not mathematical feasibility; CasADi time is not full-stack deployment latency \\
    Zero observed binary failures & Native callbacks, zero-margin physical overlap and 0.25-m-per-actor margin-adjusted separation are reported separately & No zero-risk or equivalence claim \\
    Very small active tail & Separate labelled legacy diagnostic & 15 windows/six rollouts/three groups; no stable tail estimate \\
    Legacy implementation differs & Corrected R3 is primary; no pooling & Legacy evidence is mechanism-only \\
    \bottomrule
  \end{tabular}
\end{table}

The fixed policies are three prespecified operating points rather than a dense
continuous Pareto frontier. The implementation diagram represents information
flow, not a causal graph. These boundaries prevent overclaiming but do not
invalidate the frozen in-distribution and corrected R3 comparisons.

\subsection{Implications and future work}

For research practice, a simple task-adaptation control and training-budget
audit should precede claims for a more elaborate interaction architecture.
Closed-loop evaluation should verify that the prediction manipulation persists
in deployment and should retain policy/context cells rather than collapse them
into one score. For system design, continuous geometry margins and mechanism
telemetry remain informative when nominal binary failure counts are zero.

Future work should replicate across maps and junctions, increase independent
initialisation clusters, enrich response-active data, compare calibration-
matched predictor stacks, and evaluate larger graph/attention models with
training budgets commensurate with their capacity. These are external-validity
extensions, not missing experiments required to support the present claims.

For production-oriented evaluation, the immediate implication is procedural:
retain a cheap task-adapted control, audit probabilistic calibration by
interaction subgroup, verify that a model difference survives in-loop, and
report a policy frontier with execution mechanisms. The present sample cannot
set deployment thresholds, but this sequence reduces the risk of promoting a
larger model or an adaptive policy on an offline score alone.
